\documentclass[12pt]{article}
\usepackage[utf8]{inputenc}
\usepackage[T1]{fontenc}
\usepackage[margin=1in]{geometry}
\usepackage{setspace}
\usepackage{xcolor}
\usepackage[colorlinks=true, urlcolor=blue!60!black, linkcolor=blue!60!black]{hyperref}
\setlength{\parindent}{0em}
\setlength{\parskip}{0.45em}
\renewcommand{\baselinestretch}{1.0}
\begin{document}
\begin{center}
{\Large \textit{Katzenbach v. Morgan}, 384 U.S. 641 (1966)\par}
\vspace{0.3em}
{\normalsize Case Brief --- One-Page Edition\par}
\end{center}

\textbf{Memory Jogger.} Section 5 lets Congress override New York's English literacy test for Puerto Rican voters without a prior judicial ruling (\href{https://supreme.justia.com/cases/federal/us/384/641/}{Justia}).

\textbf{Facts.} New York's English literacy requirement disenfranchised many Spanish-schooled Puerto Rican migrants. Section 4(e) of the Voting Rights Act of 1965 protected anyone with a sixth-grade education in an ``American-flag'' school in Puerto Rico. City voters challenged it.

\textbf{Procedural History.} A three-judge District Court for the District of Columbia held 2--1 that Congress exceeded its power and enjoined \S 4(e).

\textbf{Votes.} Reversed, 7--2 (Brennan for the Court; Douglas in part; Harlan, with Stewart, dissenting).

\textbf{Holding.} Section 4(e) is appropriate \S 5 legislation preempting New York's test; Congress need not await a judicial finding, and it is enough that the Court can perceive a rational basis for its judgment.

\textbf{Reasoning.} A judicial-finding rule would reduce Congress to ``the insignificant role of abrogating only those state laws that the judicial branch was prepared to adjudge unconstitutional.'' Congress could find the vote would protect Puerto Ricans in public services, or itself judge the test invidious; Footnote 10 forbids diluting rights. \emph{Separately:} Harlan called the question ``a judicial question'' and noted Congress built no record. \textit{City of Boerne} later kept only the remedial rationale.

\textbf{Examples.} \emph{Same-Side:} a nationwide literacy-test ban. \emph{Opposite-Side:} RFRA's override of a Free Exercise ruling. \emph{Fence-Sitter:} an 18-year-old vote in state elections.

\textbf{Critique.} Faulted for letting Congress redefine rights; defended as democratic constitutionalism.

\textbf{Key Quotation.} ``Thus, the \textit{McCulloch v. Maryland} standard is the measure of what constitutes `appropriate legislation' under \S 5 of the Fourteenth Amendment. Correctly viewed, \S 5 is a positive grant of legislative power authorizing Congress to exercise its discretion in determining whether and what legislation is needed to secure the guarantees of the Fourteenth Amendment.'' --- Justice Brennan, 384 U.S. at 651.

\end{document}
